\begin{proof}[Proof of \cref{obj:thm:annotation-resource-phases}]
\begin{enumerate}
\item \emph{Uniform risk comparison.}
By \cref{obj:thm:uniform-annotation-frontier}, there are numerical constants \(0<c_{\mathrm F}\le C_{\mathrm F}<\infty\) such that, throughout the allowed public index range,
\[
c_{\mathrm F}r(n,m,d,\epsilon)
\le R(n,m,d,\epsilon)
\le C_{\mathrm F}r(n,m,d,\epsilon),
\]
where \(r\) is defined in \cref{obj:def:rate-handle}.

Fix an allowed sequence \(\mathbf v\), and write
\[
\begin{gathered}
S_k=n_k\epsilon_k,\qquad
N_k=n_k+m_k,\qquad
\ell_k=\log(\exp(1)+S_k),\\
R_k=R(n_k,m_k,d_k,\epsilon_k),\qquad
r_k=r(n_k,m_k,d_k,\epsilon_k),\\
b_k=b(n_k,\epsilon_k)=\min\{1,S_k^{-1}\},\qquad
z_k=\frac{d_k}{N_k\epsilon_k\ell_k}.
\end{gathered}
\]
The public restrictions give \(S_k>0\), \(N_k>0\), and \(\ell_k>0\). Consequently,
\[
z_k>0,\qquad
r_k=\min\{1,S_k^{-1}+z_k^2\}>0,\qquad
0<c_{\mathrm F}r_k\le R_k\le C_{\mathrm F}r_k.
\]
The same comparison and positivity apply at every auxiliary budget, including zero.
% lean: CausalSmith.Stat.AnnotationRarearmFrontier.annotation_resource_phases

\item \emph{Consistency.}
The comparison gives
\[
0\le r_k\le \frac{R_k}{c_{\mathrm F}},
\qquad
0\le R_k\le C_{\mathrm F}r_k,
\]
so \(R_k\to0\) if and only if \(r_k\to0\). If \(r_k\to0\), then eventually \(r_k<1\), and hence
\[
r_k=S_k^{-1}+z_k^2.
\]
Both summands are nonnegative, so
\[
0<S_k^{-1}\le r_k\longrightarrow0,
\qquad
0\le z_k^2\le r_k\longrightarrow0.
\]
Taking reciprocals in the first limit and nonnegative square roots in the second yields \(S_k\to\infty\) and \(z_k\to0\). Conversely, these two limits imply
\[
0\le r_k\le S_k^{-1}+z_k^2\longrightarrow0.
\]
Since \(N_k\epsilon_k\ell_k>0\),
\[
z_k\longrightarrow0
\quad\Longleftrightarrow\quad
d_k=o(N_k\epsilon_k\ell_k).
\]
Together with \cref{obj:def:phases}, this proves the consistency characterization.
% lean: CausalSmith.Stat.AnnotationRarearmFrontier.consistency_of_frontier_comparison

\item \emph{Oracle membership.}
Suppose \(S_k\to\infty\). Then eventually \(S_k\ge1\), and
\[
b_k=S_k^{-1}.
\]
By \cref{obj:def:phases}, the remaining oracle condition is eventual boundedness above of \(R_k/b_k\).

First suppose that, for some real constant \(M_{\mathrm{or}}\),
\[
\frac{R_k}{b_k}\le M_{\mathrm{or}}
\]
eventually. Define
\[
K_{\mathrm{or}}=\frac{\max\{M_{\mathrm{or}},1\}}{c_{\mathrm F}}>0.
\]
The lower comparison then gives, eventually,
\[
c_{\mathrm F}r_k\le R_k
\le \frac{\max\{M_{\mathrm{or}},1\}}{S_k},
\qquad
r_k\le\frac{K_{\mathrm{or}}}{S_k}.
\]
Eventually \(S_k>K_{\mathrm{or}}\), so \(r_k<1\) and its cap is inactive. Thus
\[
1+S_kz_k^2=S_kr_k\le K_{\mathrm{or}},
\qquad
(z_k\sqrt{S_k})^2\le K_{\mathrm{or}}.
\]
All factors are nonnegative, and therefore
\[
z_k\sqrt{S_k}\le\sqrt{K_{\mathrm{or}}},
\qquad
d_k\le
\sqrt{K_{\mathrm{or}}}\,
\frac{N_k\epsilon_k\ell_k}{\sqrt{S_k}}
\]
eventually. This is the required big-\(O\) condition.

Conversely, that big-\(O\) condition supplies a constant \(K_{\mathrm{dim}}>0\) such that eventually
\[
d_k\le K_{\mathrm{dim}}
\frac{N_k\epsilon_k\ell_k}{\sqrt{S_k}}.
\]
Multiplication by the positive factor
\(\sqrt{S_k}/(N_k\epsilon_k\ell_k)\) gives
\[
z_k\sqrt{S_k}\le K_{\mathrm{dim}},
\qquad
S_kz_k^2\le K_{\mathrm{dim}}^2.
\]
Consequently,
\[
r_k\le S_k^{-1}+z_k^2
\le\frac{1+K_{\mathrm{dim}}^2}{S_k},
\qquad
\frac{R_k}{b_k}
\le C_{\mathrm F}(1+K_{\mathrm{dim}}^2)
\]
eventually. This proves the oracle equivalence.
% lean: CausalSmith.Stat.AnnotationRarearmFrontier.oracle_of_frontier_comparison

\item \emph{Necessary conditions for improvement.}
Define the supervised risk, supervised rate, and rate ratio by
\[
R_{\mathrm{sup},k}=R(n_k,0,d_k,\epsilon_k),\qquad
r_{\mathrm{sup},k}=r(n_k,0,d_k,\epsilon_k),\qquad
\rho_k=\frac{r_k}{r_{\mathrm{sup},k}}.
\]
Their denominators are positive. Applying the uniform comparison at both auxiliary budgets gives
\[
\frac{c_{\mathrm F}}{C_{\mathrm F}}\rho_k
\le \frac{R_k}{R_{\mathrm{sup},k}}
\le \frac{C_{\mathrm F}}{c_{\mathrm F}}\rho_k.
\]
Indeed, the left inequality follows by dividing
\(R_k\ge c_{\mathrm F}r_k\) by
\(R_{\mathrm{sup},k}\le C_{\mathrm F}r_{\mathrm{sup},k}\);
the right inequality follows from the opposite two bounds. Squeezing therefore yields
\[
\mathbf v\in\mathcal I
\quad\Longleftrightarrow\quad
\rho_k\longrightarrow0.
\]
% lean: CausalSmith.Stat.AnnotationRarearmFrontier.improvement_ratio_of_frontier_comparison

Assume \(\rho_k\to0\). Since \(0<r_{\mathrm{sup},k}\le1\),
\[
0\le r_k=\rho_kr_{\mathrm{sup},k}\le\rho_k\longrightarrow0.
\]
The consistency calculation gives \(S_k\to\infty\).
% lean: CausalSmith.Stat.AnnotationRarearmFrontier.phase_improvement_label_divergence

Introduce, for every \(k\),
\[
\xi_k=\frac{d_k}{S_k\ell_k}>0,\qquad
\eta_k=\frac{N_k}{n_k}>0,\qquad
a_{\mathrm{lab},k}=S_k^{-1}>0,\qquad
t_{\mathrm{cap},k}=\min\{1,\xi_k^2\}>0.
\]
The identity \(N_k\epsilon_k=\eta_kS_k\) gives
\[
r_k=\min\{1,a_{\mathrm{lab},k}+(\xi_k/\eta_k)^2\},
\qquad
r_{\mathrm{sup},k}=\min\{1,a_{\mathrm{lab},k}+\xi_k^2\}.
\]

Eventually \(S_k\ge1\) and \(0\le\rho_k\le1/2\). For such \(k\),
\[
a_{\mathrm{lab},k}\le1,\qquad
r_k=\rho_kr_{\mathrm{sup},k}\le\rho_k\le\frac12.
\]
Hence the cap in \(r_k\) is inactive:
\[
r_k=a_{\mathrm{lab},k}+(\xi_k/\eta_k)^2.
\]
In particular,
\[
a_{\mathrm{lab},k}
\le r_k
=\rho_kr_{\mathrm{sup},k}
\le\rho_k(a_{\mathrm{lab},k}+\xi_k^2)
\le\frac12(a_{\mathrm{lab},k}+\xi_k^2).
\]
It follows that
\[
a_{\mathrm{lab},k}\le\xi_k^2,
\qquad
\frac12a_{\mathrm{lab},k}
\le(1-\rho_k)a_{\mathrm{lab},k}
\le\rho_k\xi_k^2.
\]
Using \((\sqrt{S_k})^2=S_k\), we obtain the first inverse-square bound:
\[
(\sqrt{S_k}\xi_k)^{-2}
=\frac{a_{\mathrm{lab},k}}{\xi_k^2}
\le2\rho_k.
\]

For the second bound, the cap satisfies, for every \(k\),
\[
t_{\mathrm{cap},k}\le r_{\mathrm{sup},k}
\le t_{\mathrm{cap},k}+a_{\mathrm{lab},k}.
\]
The lower bound follows by increasing the second argument of the minimum. For the upper bound, if \(\xi_k^2\le1\), then
\[
r_{\mathrm{sup},k}\le a_{\mathrm{lab},k}+\xi_k^2
=a_{\mathrm{lab},k}+t_{\mathrm{cap},k};
\]
if \(\xi_k^2>1\), then
\[
r_{\mathrm{sup},k}\le1=t_{\mathrm{cap},k}
\le t_{\mathrm{cap},k}+a_{\mathrm{lab},k}.
\]
% lean: CausalSmith.Stat.AnnotationRarearmFrontier.phase_supervised_cap_bounds

For the eventual indices under consideration,
\(a_{\mathrm{lab},k}\le1\) and
\(a_{\mathrm{lab},k}\le\xi_k^2\), so
\[
a_{\mathrm{lab},k}\le t_{\mathrm{cap},k},
\qquad
r_{\mathrm{sup},k}\le2t_{\mathrm{cap},k}.
\]
Thus
\[
(\xi_k/\eta_k)^2
\le r_k
=\rho_kr_{\mathrm{sup},k}
\le2\rho_kt_{\mathrm{cap},k}.
\]
The exact budget identity is
\[
\frac{(\xi_k/\eta_k)^2}{t_{\mathrm{cap},k}}
=
\left(\frac{\max\{1,\xi_k\}}{\eta_k}\right)^2
=
\left(\frac{\eta_k}{\max\{1,\xi_k\}}\right)^{-2}.
\]
If \(\xi_k\le1\), the first two expressions both equal
\(\eta_k^{-2}\); if \(\xi_k>1\), they both equal
\((\xi_k/\eta_k)^2\). Therefore,
\[
\left(\frac{\eta_k}{\max\{1,\xi_k\}}\right)^{-2}
\le2\rho_k.
\]
% lean: CausalSmith.Stat.AnnotationRarearmFrontier.phase_improvement_necessary_bounds

Both inverse squares tend to zero. Taking positive square roots and then reciprocals gives
\[
\sqrt{S_k}\xi_k\longrightarrow\infty,
\qquad
\frac{\eta_k}{\max\{1,\xi_k\}}\longrightarrow\infty.
\]
Finally,
\[
\sqrt{S_k}\xi_k
=\frac{d_k}{\sqrt{S_k}\ell_k},
\qquad
\frac{\eta_k}{\max\{1,\xi_k\}}
=
\frac{N_k/n_k}{\max\{1,d_k/(S_k\ell_k)\}},
\]
which proves all three necessary limits.
% lean: CausalSmith.Stat.AnnotationRarearmFrontier.phase_improvement_ratio_necessary

\item \emph{Sufficient conditions for improvement.}
Suppose the three limits in the claimed improvement characterization hold. In the notation just defined, they are
\[
S_k\longrightarrow\infty,\qquad
\sqrt{S_k}\xi_k\longrightarrow\infty,\qquad
\frac{\eta_k}{\max\{1,\xi_k\}}\longrightarrow\infty.
\]
For every \(k\), positivity and
\(t_{\mathrm{cap},k}\le r_{\mathrm{sup},k}\) give
\[
0\le\rho_k
\le
\frac{a_{\mathrm{lab},k}+(\xi_k/\eta_k)^2}
     {t_{\mathrm{cap},k}}
=
\frac{a_{\mathrm{lab},k}}{t_{\mathrm{cap},k}}
+
\frac{(\xi_k/\eta_k)^2}{t_{\mathrm{cap},k}}.
\]
The first quotient satisfies
\[
\frac{a_{\mathrm{lab},k}}{t_{\mathrm{cap},k}}
=
\begin{cases}
(\sqrt{S_k}\xi_k)^{-2},&\xi_k^2\le1,\\
S_k^{-1},&\xi_k^2>1,
\end{cases}
\le S_k^{-1}+(\sqrt{S_k}\xi_k)^{-2}.
\]
Combining this with the budget identity yields
\[
0\le\rho_k
\le
S_k^{-1}
+(\sqrt{S_k}\xi_k)^{-2}
+\left(\frac{\eta_k}{\max\{1,\xi_k\}}\right)^{-2}
\longrightarrow0.
\]
The risk-ratio comparison then proves
\(\mathbf v\in\mathcal I\).
% lean: CausalSmith.Stat.AnnotationRarearmFrontier.phase_improvement_ratio_sufficient

\item \emph{The supervised slice.}
If \(m_k=0\) for every \(k\), then
\[
N_k\epsilon_k=S_k,\qquad
\frac{N_k\epsilon_k\ell_k}{\sqrt{S_k}}
=\frac{S_k\ell_k}{\sqrt{S_k}}
=\sqrt{S_k}\ell_k,
\]
where the final equality uses \(S_k=(\sqrt{S_k})^2\) and
\(\sqrt{S_k}>0\). Substitution into the consistency and oracle characterizations gives the two stated supervised conditions.

Moreover, \(R_{\mathrm{sup},k}=R_k>0\), so
\[
\frac{R_k}{R_{\mathrm{sup},k}}=1
\qquad\text{for every }k.
\]
Its limit is \(1\), and hence \(\mathbf v\notin\mathcal I\).
% lean: CausalSmith.Stat.AnnotationRarearmFrontier.annotation_resource_phases

\item \emph{The equivalent oracle budget.}
For every \(k\) and every constant \(K_{\mathrm{bud}}\ge0\), multiplication or division by the positive quantities
\(\sqrt{S_k}\) and \(\epsilon_k\ell_k\) gives
\[
d_k\le K_{\mathrm{bud}}
\frac{N_k\epsilon_k\ell_k}{\sqrt{S_k}}
\quad\Longleftrightarrow\quad
\frac{d_k\sqrt{S_k}}{\epsilon_k\ell_k}
\le K_{\mathrm{bud}}N_k.
\]
Thus the two eventual inequalities use the same comparison constant, and
\[
d_k=O\!\left(\frac{N_k\epsilon_k\ell_k}{\sqrt{S_k}}\right)
\quad\Longleftrightarrow\quad
\frac{d_k\sqrt{S_k}}{\epsilon_k\ell_k}=O(N_k).
\]
Under \(S_k\to\infty\), the oracle characterization proves the asserted budget equivalence.
% lean: CausalSmith.Stat.AnnotationRarearmFrontier.phase_dimension_budget_iff

\item \emph{Risk order and the uniform auxiliary-budget floor.}
For any allowed experiment, increasing the second argument of the minimum gives
\[
b(n,\epsilon)
=\min\{1,(n\epsilon)^{-1}\}
\le
\min\left\{1,(n\epsilon)^{-1}
+\left(\frac{d}{(n+m)\epsilon\ell}\right)^2\right\}
=r(n,m,d,\epsilon).
\]
Hence, for every \(k\) and every \(m'\in\mathbb N\),
\[
c_{\mathrm F}b_k
\le c_{\mathrm F}r(n_k,m',d_k,\epsilon_k)
\le R(n_k,m',d_k,\epsilon_k).
\]
% lean: CausalSmith.Stat.AnnotationRarearmFrontier.phase_labelBenchmark_le_frontier

Now let \(\mathbf v\in\mathcal O\). Choose a real constant
\(M_{\mathrm{risk}}\) giving an eventual upper bound for \(R_k/b_k\), and define
\[
C_{\mathrm{or}}=\max\{M_{\mathrm{risk}},1\}>0.
\]
Since \(S_k\to\infty\), eventually \(b_k=S_k^{-1}\). On that same tail,
\[
\frac{c_{\mathrm F}}{S_k}\le R_k
\le\frac{M_{\mathrm{risk}}}{S_k}
\le\frac{C_{\mathrm{or}}}{S_k}.
\]
The lower comparison also gives, on a tail independent of \(m'\),
\[
\frac{c_{\mathrm F}}{S_k}
\le R(n_k,m',d_k,\epsilon_k)
\qquad\text{for every }m'\in\mathbb N.
\]
This proves both oracle risk assertions with a positive lower constant independent of the auxiliary budget.
% lean: CausalSmith.Stat.AnnotationRarearmFrontier.annotation_resource_phases

\item \emph{Single-experiment saturation.}
Fix an experiment satisfying the stated public restrictions, and use
\[
N=n+m,\qquad
S=n\epsilon,\qquad
\ell=\log(\exp(1)+S).
\]
All three denominator factors are positive. If \(N\epsilon\ell\le d\), then
\[
\frac{d}{N\epsilon\ell}\ge1,
\qquad
S^{-1}+\left(\frac{d}{N\epsilon\ell}\right)^2\ge1,
\qquad
r(n,m,d,\epsilon)=1.
\]
The uniform comparison therefore yields
\[
c_{\mathrm F}\le R(n,m,d,\epsilon).
\]
% lean: CausalSmith.Stat.AnnotationRarearmFrontier.phase_frontier_saturation

For the upper bound, the laws in
\cref{obj:thm:rare-label-floor} ensure that the model class is nonempty. For each law \(P\) in that class, the binary outcome conditional means lie in \([0,1]\). The target formula in \cref{obj:def:target} thus gives
\[
|\tau(P)|
\le\sum_{j=1}^{d}p_j(P)
       |\mu_{1j}(P)-\mu_{0j}(P)|
\le\sum_{j=1}^{d}p_j(P)=1.
\]
Define the constant original-record rule by
\[
T_{\mathrm{zero}}(\mathcal D,U)=0
\]
for every data realization and seed. This rule is Borel measurable and takes values in \([-1,1]\). Its squared error is the constant \(\tau(P)^2\), so \cref{obj:def:risk} gives
\[
R(n,m,d,\epsilon)
\le
\sup_{P\in\mathcal M_{d,\epsilon}}
\mathbb E_P\!\left[
(T_{\mathrm{zero}}(\mathcal D,U)-\tau(P))^2
\right]
=
\sup_{P\in\mathcal M_{d,\epsilon}}\tau(P)^2
\le1.
\]
This proves the saturation bounds with the universal constant \(c_{\mathrm F}\).
% lean: CausalSmith.Stat.AnnotationRarearmFrontier.phase_minimaxRisk_le_one

\item \emph{Single-experiment oracle plateau.}
Suppose
\[
\frac{d\sqrt S}{\epsilon\ell}\le N.
\]
Multiplication by \(\epsilon\ell>0\) and division by
\(N\epsilon\ell>0\) give
\[
d\sqrt S\le N\epsilon\ell,
\qquad
\frac{d}{N\epsilon\ell}\sqrt S\le1.
\]
Squaring, using nonnegativity and \((\sqrt S)^2=S\), and dividing by \(S>0\), we obtain
\[
\left(\frac{d}{N\epsilon\ell}\right)^2\le S^{-1}.
\]
If \(S^{-1}\le1\), then
\[
r(n,m,d,\epsilon)
\le S^{-1}+\left(\frac{d}{N\epsilon\ell}\right)^2
\le2S^{-1}
=2b(n,\epsilon).
\]
If \(S^{-1}>1\), then
\[
b(n,\epsilon)=1,\qquad
r(n,m,d,\epsilon)\le1\le2b(n,\epsilon).
\]
Together with the universal lower comparison
\(b(n,\epsilon)\le r(n,m,d,\epsilon)\), these cases give
\[
c_{\mathrm F}b(n,\epsilon)
\le R(n,m,d,\epsilon)
\le C_{\mathrm F}r(n,m,d,\epsilon)
\le2C_{\mathrm F}b(n,\epsilon).
\]
The constants \(c_{\mathrm F}\) and \(2C_{\mathrm F}\) are positive and independent of all public indices.
% lean: CausalSmith.Stat.AnnotationRarearmFrontier.phase_frontier_plateau
\end{enumerate}
\end{proof}
